\section{Definitions as Machines}

Definitions are the machinery of proof. A definition is not decoration around a theorem. It is not a vocabulary word you memorize so that you can understand the sentence. In proof-based mathematics, a definition is an exact little machine with two doors.

One door is the \emph{use} door. If you are told that an object satisfies a definition, you may open the definition and take out the data it promises.

The other door is the \emph{prove} door. If you want to show that an object satisfies a definition, the definition tells you exactly what data you must build.

For example, the definition of even integer is:
\[
n\text{ is even}\quad\Longleftrightarrow\quad \exists k\in\Z\text{ such that }n=2k.
\]
If someone gives you ``$n$ is even,'' you may write $n=2k$ for some integer $k$. If someone asks you to prove ``$n$ is even,'' you must produce an integer $k$ such that $n=2k$.

\begin{focusbox}
A definition is a reversible contract: it says what you get to use and what you must prove.
\end{focusbox}

This is a deep shift from calculation-based mathematics. In a calculation course, a phrase like ``continuous function'' may often mean ``a function whose graph has no breaks.'' That picture is useful. But in proof, the picture is not the machine. The machine is the definition. Depending on the course, continuity may mean an epsilon-delta condition, preservation of limits, inverse images of open sets, or something else equivalent in a particular setting. The proof has to use the machine that has actually been installed.

\subsection{Why definitions feel weird at first}

Beginners often experience definitions as strangely artificial. Why say ``there exists an integer $k$ such that $n=2k$'' when we already know what even means? Why say ``$a$ divides $b$ if $b=ak$ for some integer $k$'' when everyone knows division?

The answer is that ordinary understanding is not stable enough for proof. It is too full of pictures, memories, and unstated assumptions. Proof needs handles.

The definition
\[
n=2k\quad\text{for some }k\in\Z
\]
gives us a handle. It turns the word ``even'' into an equation. Once we have an equation, we can substitute, add, multiply, factor, and compare. A vague feeling becomes an object we can use.

\begin{intuitionbox}
Definitions are adapters. They convert ordinary mathematical words into proof-compatible data.
\end{intuitionbox}

The same thing happens in software. A type signature does not tell you the full personality of a function, but it tells you how the function can be connected to other code. A function declared as
\[
\texttt{isEven : Integer -> Boolean}
\]
is not the same kind of thing as a function declared as
\[
\texttt{halve : EvenInteger -> Integer}.
\]
The contract matters. A compiler is strict because ambiguity makes large systems fragile. Mathematics is strict for the same reason.

\subsection{Opening a definition}

To open a definition is to replace a word by the data the definition guarantees.

Suppose a proof begins:

\begin{quote}
Let $n$ be an even integer.
\end{quote}

The next legal move is:
\[
n=2k\text{ for some }k\in\Z.
\]

Not because we guessed it. Not because it is usually true. Because that is exactly what ``even'' means.

Here is the same idea with divisibility.

\[
a\mid b
\]
is read ``$a$ divides $b$.'' The definition is:
\[
a\mid b\quad\Longleftrightarrow\quad \exists k\in\Z\text{ such that }b=ak.
\]

So if a proof says $a\mid b$, you may introduce an integer $k$ with $b=ak$.

\begin{formalbox}
Opening patterns:
\[
\begin{array}{rcl}
n\text{ is even} &\Rightarrow& n=2k\text{ for some }k\in\Z,\\[0.3em]
n\text{ is odd} &\Rightarrow& n=2k+1\text{ for some }k\in\Z,\\[0.3em]
a\mid b &\Rightarrow& b=ak\text{ for some }k\in\Z,\\[0.3em]
x\in A\cap B &\Rightarrow& x\in A\text{ and }x\in B,\\[0.3em]
x\in A\cup B &\Rightarrow& x\in A\text{ or }x\in B.
\end{array}
\]
\end{formalbox}

Opening a definition is often the first real move in a proof. A theorem gives you hypotheses. The hypotheses are wrapped packages. Definitions unwrap them.

\subsection{Closing a definition}

Closing a definition goes in the other direction. You have built the data, and now you are allowed to name the property.

Suppose you have reached
\[
m+n=2(r+s)
\]
where $r+s\in\Z$. Then you may conclude $m+n$ is even. You are closing the definition of even.

Suppose you have reached
\[
c=a(mn)
\]
where $mn\in\Z$. Then you may conclude $a\mid c$. You are closing the definition of divisibility.

\begin{focusbox}
Opening turns a word into usable data. Closing turns constructed data into the word you were trying to prove.
\end{focusbox}

This open-close rhythm is everywhere:

\begin{itemize}
\item To use injectivity, open it: from $f(x)=f(y)$ conclude $x=y$.
\item To prove injectivity, close it: start with arbitrary $x,y$ and show $f(x)=f(y)\Rightarrow x=y$.
\item To use subset inclusion, open it: from $A\subseteq B$ and $x\in A$ conclude $x\in B$.
\item To prove subset inclusion, close it: let $x\in A$ be arbitrary and show $x\in B$.
\end{itemize}

The words change. The machine rhythm stays the same.

\subsection{Worked proof: even plus even}

\begin{theorem}
The sum of two even integers is even.
\end{theorem}

\begin{proof}
Let $a,b\in\Z$ be even. Since $a$ is even, there exists $m\in\Z$ such that $a=2m$. Since $b$ is even, there exists $n\in\Z$ such that $b=2n$. Then
\[
a+b=2m+2n=2(m+n).
\]
Because $m+n\in\Z$, this shows that $a+b$ is twice an integer. Therefore $a+b$ is even.
\end{proof}

This proof is small, but it contains the whole pattern.

We opened the definition twice:
\[
a=2m,\qquad b=2n.
\]
We used algebra:
\[
a+b=2(m+n).
\]
Then we closed the definition:
\[
a+b\text{ is even}.
\]

The proof is not proving that $2m+2n=2(m+n)$ in isolation. That algebra line is just the hinge. The real proof structure is definition, algebra, definition.

\begin{walkbox}
When you read a proof, ask: which definition was opened here, and which definition was closed at the end?
\end{walkbox}

\subsection{The integer witness matters}

The phrase ``for some integer $k$'' is not filler. It is part of the machine.

To prove $n$ is even, it is not enough to write
\[
n=2\cdot\frac n2.
\]
That is true for every real number $n$, but it does not prove that every real number is even. The witness must be an integer.

For instance, take $n=3$. Then
\[
3=2\cdot\frac32.
\]
But $\frac32\notin\Z$, so this does not show that $3$ is even.

\begin{trapbox}
Whenever a definition says ``there exists $k\in\Z$,'' the membership condition $k\in\Z$ is not optional. Many false proofs accidentally use a witness from the wrong universe.
\end{trapbox}

Here is a correct closing step:
\[
6n+10=2(3n+5).
\]
If $n\in\Z$, then $3n+5\in\Z$. Therefore $6n+10$ is even.

Here is a suspicious closing step:
\[
3n+1=2\left(\frac{3n+1}{2}\right).
\]
This does not prove $3n+1$ is even unless you already know $\frac{3n+1}{2}$ is an integer. For $n=1$, it is $2$; for $n=2$, it is $\frac72$. The witness test catches the mistake.

\subsection{Worked proof: odd plus odd}

\begin{theorem}
The sum of two odd integers is even.
\end{theorem}

\begin{proof}
Let $a,b\in\Z$ be odd. Since $a$ is odd, there exists $m\in\Z$ such that $a=2m+1$. Since $b$ is odd, there exists $n\in\Z$ such that $b=2n+1$. Then
\[
a+b=(2m+1)+(2n+1)=2m+2n+2=2(m+n+1).
\]
Since $m+n+1\in\Z$, the sum $a+b$ is twice an integer. Therefore $a+b$ is even.
\end{proof}

Notice that the proof did not say ``odd plus odd is obviously even.'' It translated oddness into equations and translated the final equation into evenness. The obvious idea was real, but the definition made it proof.

\subsection{Worked proof: divisibility is transitive}

\begin{theorem}
Let $a,b,c\in\Z$. If $a\mid b$ and $b\mid c$, then $a\mid c$.
\end{theorem}

\begin{proof}
Suppose $a\mid b$ and $b\mid c$. Since $a\mid b$, there exists $m\in\Z$ such that $b=am$. Since $b\mid c$, there exists $n\in\Z$ such that $c=bn$. Substituting $b=am$ into $c=bn$, we get
\[
c=(am)n=a(mn).
\]
Because $mn\in\Z$, this shows that $c$ is an integer multiple of $a$. Hence $a\mid c$.
\end{proof}

The proof is just the definition of divisibility doing its job. The input data gives:
\[
b=am,\qquad c=bn.
\]
The desired conclusion is $a\mid c$, so we want:
\[
c=a(\text{integer}).
\]
Substitution gives exactly that.

\begin{trapbox}
The statement $a\mid b$ means $b=ak$, not $a=bk$. The number being divided is the one that becomes a multiple.
\end{trapbox}

\subsection{Bad proof debugged: divisibility}

Consider the claim:

\begin{theorem}
If $a\mid b$ and $a\mid c$, then $a\mid(b+c)$.
\end{theorem}

Here is a bad proof:

\begin{quote}
Since $a\mid b$ and $a\mid c$, divide both by $a$. Then $b/a$ and $c/a$ are integers. So $(b+c)/a$ is an integer. Therefore $a\mid(b+c)$.
\end{quote}

The idea is not terrible, but the proof is slippery. It uses fractions where integer equations would be cleaner. It also hides the integer witnesses.

Here is the definition-machine version.

\begin{proof}
Suppose $a\mid b$ and $a\mid c$. Since $a\mid b$, there exists $m\in\Z$ such that $b=am$. Since $a\mid c$, there exists $n\in\Z$ such that $c=an$. Then
\[
b+c=am+an=a(m+n).
\]
Because $m+n\in\Z$, we have $a\mid(b+c)$.
\end{proof}

The corrected proof exposes the structure:

\begin{itemize}
\item $a\mid b$ gives one integer witness $m$.
\item $a\mid c$ gives another integer witness $n$.
\item The new witness for $a\mid(b+c)$ is $m+n$.
\end{itemize}

\begin{labbox}
Debug this proof:

\begin{quote}
If $a\mid b$, then $a\mid bc$. Since $a\mid b$, we know $a=bk$. Multiplying by $c$ gives $ac=bkc$, so $a\mid bc$.
\end{quote}

Find the exact line where the definition was opened incorrectly. Then rewrite the proof using the correct witness.
\end{labbox}

\subsection{Definitions with hidden quantifiers}

Many definitions hide quantifiers inside ordinary-sounding words.

\[
A\subseteq B
\]
means:
\[
\forall x,\ x\in A\Rightarrow x\in B.
\]

To use $A\subseteq B$, you need an element $x\in A$. Then you can conclude $x\in B$.

To prove $A\subseteq B$, you must show that every element of $A$ belongs to $B$. The usual proof pattern is:

\begin{quote}
Let $x\in A$ be arbitrary. We will show $x\in B$.
\end{quote}

That one sentence is not decorative. It is the machine for proving a universal statement.

\begin{theorem}
For all sets $A$ and $B$, $A\cap B\subseteq A$.
\end{theorem}

\begin{proof}
Let $x\in A\cap B$ be arbitrary. By definition of intersection, $x\in A$ and $x\in B$. In particular, $x\in A$. Therefore $A\cap B\subseteq A$.
\end{proof}

Again the proof is open-close:

\begin{itemize}
\item To prove a subset claim, open the subset target: let $x$ be an arbitrary element of the left side.
\item Open the definition of intersection.
\item Close the subset claim by showing $x$ belongs to the right side.
\end{itemize}

\subsection{Definitions in analysis: bounded sequences}

Let us look ahead to real analysis.

\begin{definition}
A sequence $(a_n)$ of real numbers is bounded above if there exists $M\in\R$ such that $a_n\leq M$ for every $n\in\N$.
\end{definition}

This definition has two important pieces:

\begin{itemize}
\item there exists one real number $M$,
\item that same $M$ works for every index $n$.
\end{itemize}

To use bounded-above-ness, you may introduce a real number $M$ and say $a_n\leq M$ for all $n$.

To prove bounded-above-ness, you must produce one number $M$ that works for all terms.

\begin{theorem}
The sequence $a_n=1-\frac1n$ is bounded above.
\end{theorem}

\begin{proof}
For every $n\in\N$, we have $\frac1n>0$, so $1-\frac1n<1$. Thus $a_n\leq1$ for every $n\in\N$. Therefore the sequence is bounded above, with upper bound $M=1$.
\end{proof}

The proof closes the definition by presenting a witness: $M=1$.

Could we choose $M=100$? Yes. A bound does not have to be the best bound. To prove bounded above, one working ceiling is enough.

Could we choose $M=1-\frac1n$? No. That expression depends on $n$. The definition asks for one $M$ that works for every $n$.

\begin{trapbox}
When a definition says $\exists M\ \forall n$, the $M$ must be chosen before $n$. It cannot secretly change from term to term.
\end{trapbox}

\subsection{Equivalent definitions and why they are allowed}

Sometimes a concept has multiple definitions that turn out to be equivalent. This can be confusing: if definitions are supposed to be exact machines, how can there be more than one?

The answer is that a mathematical theory may prove two machines always accept the same inputs. Once that theorem is established, either machine can be used.

For example, for an integer $n$, these are equivalent:

\[
n\text{ is even}
\]
and
\[
n^2\text{ is even}.
\]

But we should not silently use the second as the definition unless it has been stated or proved. The usual definition of even is $n=2k$ for some integer $k$. The statement ``$n^2$ is even iff $n$ is even'' is a theorem.

Another example comes from continuity. In a first analysis course, continuity at a point often means:
\[
\forall\eps>0,\exists\delta>0,\forall x,\ |x-a|<\delta\Rightarrow |f(x)-f(a)|<\eps.
\]
In topology, continuity is often defined by:
\[
U\text{ open in the codomain}\Rightarrow f^{-1}(U)\text{ open in the domain}.
\]
For functions $\R\to\R$ with the usual topology, these definitions agree. But proving that agreement is mathematics, not word magic.

\begin{formalbox}
Use the definition your context gives you. Use an equivalent characterization only after it has been proved, cited, or explicitly allowed.
\end{formalbox}

\subsection{Bad proof debugged: rational numbers}

\begin{definition}
A real number $r$ is rational if there exist integers $p,q$ with $q\neq0$ such that
\[
r=\frac pq.
\]
\end{definition}

Claim:

\begin{theorem}
The sum of two rational numbers is rational.
\end{theorem}

Bad proof:

\begin{quote}
Let $a$ and $b$ be rational. Then $a=p/q$ and $b=p/q$. So $a+b=2p/q$, which is rational.
\end{quote}

The proof accidentally uses the same numerator and denominator for both rational numbers. The definition gives each rational number its own witnesses. There is no reason the witnesses must match.

Correct proof:

\begin{proof}
Let $a,b\in\R$ be rational. Since $a$ is rational, there exist integers $p,q$ with $q\neq0$ such that $a=p/q$. Since $b$ is rational, there exist integers $r,s$ with $s\neq0$ such that $b=r/s$. Then
\[
a+b=\frac pq+\frac rs=\frac{ps+rq}{qs}.
\]
Now $ps+rq\in\Z$, $qs\in\Z$, and $qs\neq0$. Therefore $a+b$ is rational.
\end{proof}

The repaired proof respects the witnesses:

\[
a=\frac pq,\qquad b=\frac rs.
\]

Different objects usually need different witness names. Reusing letters is not illegal when done carefully, but it is a common way to lie to yourself.

\begin{trapbox}
If two hypotheses both say ``there exists,'' do not assume they give the same object. Give the witnesses different names unless there is a reason they must coincide.
\end{trapbox}

\subsection{Definitions as tests}

Definitions also tell you how to find examples and non-examples.

To test whether $18$ is divisible by $6$, ask whether
\[
18=6k
\]
for some integer $k$. Yes: $k=3$.

To test whether $20$ is divisible by $6$, ask whether
\[
20=6k
\]
for some integer $k$. No integer $k$ works.

To test whether the set
\[
A=\{2,4,6\}
\]
is a subset of
\[
B=\{1,2,3,4,5,6\},
\]
check every element of $A$. Yes, each belongs to $B$.

To test whether $B\subseteq A$, find one element of $B$ not in $A$. The number $1$ works, so $B\not\subseteq A$.

The definition tells you what evidence counts.

\subsection{A larger worked proof: product of odds}

\begin{theorem}
The product of two odd integers is odd.
\end{theorem}

\begin{proof}
Let $a,b\in\Z$ be odd. Then there exist $m,n\in\Z$ such that
\[
a=2m+1,\qquad b=2n+1.
\]
Thus
\[
ab=(2m+1)(2n+1)=4mn+2m+2n+1.
\]
Factor out $2$ from all terms except the final $1$:
\[
ab=2(2mn+m+n)+1.
\]
Since $2mn+m+n\in\Z$, the product $ab$ has the form $2k+1$ for an integer $k$. Therefore $ab$ is odd.
\end{proof}

The witness for the final oddness is
\[
k=2mn+m+n.
\]
The proof is easier to trust when you can point to the witness.

\subsection{Walking questions}

\begin{walkbox}
Pick a definition you know: prime, continuous, invertible, independent, open, closed, differentiable. Ask:

\begin{enumerate}
\item What data do I get if I am told the definition holds?
\item What data must I build if I want to prove it holds?
\item Which words in the definition are easy to ignore but actually essential?
\end{enumerate}
\end{walkbox}

\section{Proof Anatomy and Proof Compression}

Real mathematical writing is compressed. A textbook proof may use three lines where a beginner needs a page. A research paper may use one phrase where a textbook uses a paragraph. This is not because mathematicians enjoy hiding the ball. It is because communities develop shared proof idioms. To read serious mathematics, you need to learn how to expand compressed proofs and how to compress your own proofs responsibly.

\subsection{Proof anatomy}

Every proof has anatomy:

\begin{enumerate}
\item \textbf{Claim:} What exactly is being proved?
\item \textbf{Assumptions:} What are you allowed to use?
\item \textbf{Target:} What must be shown by definition?
\item \textbf{Available tools:} Which definitions, previous theorems, and algebraic facts are relevant?
\item \textbf{Strategy:} Direct, contrapositive, contradiction, cases, induction, construction, or two inclusions?
\item \textbf{Intermediate states:} What smaller claims bridge assumptions to target?
\item \textbf{Closure:} Where exactly does the proof reach the target?
\end{enumerate}

\begin{center}
\begin{tikzpicture}[
  node distance=0.9cm,
  box/.style={draw, rounded corners, align=center, minimum width=3.7cm, minimum height=0.75cm, fill=gray!5},
  arrow/.style={-{Latex}, thick}
]
\node[box] (claim) {Claim};
\node[box, below=of claim] (parse) {Parse assumptions\\and target};
\node[box, below=of parse] (tools) {Open definitions\\choose tools};
\node[box, below=of tools] (strategy) {Choose strategy};
\node[box, below=of strategy] (steps) {Build intermediate\\states};
\node[box, below=of steps] (close) {Close target\\by definition};
\draw[arrow] (claim) -- (parse);
\draw[arrow] (parse) -- (tools);
\draw[arrow] (tools) -- (strategy);
\draw[arrow] (strategy) -- (steps);
\draw[arrow] (steps) -- (close);
\end{tikzpicture}
\end{center}

\subsection{Anatomy example: composition of injective functions}

Claim:

\[
\text{If }f:A\to B\text{ and }g:B\to C\text{ are injective, then }g\circ f\text{ is injective.}
\]

\textbf{Assumptions.}

\begin{itemize}
\item $f:A\to B$.
\item $g:B\to C$.
\item $f$ is injective.
\item $g$ is injective.
\end{itemize}

\textbf{Target.}

To prove $g\circ f$ is injective, we must show:
\[
(g\circ f)(a_1)=(g\circ f)(a_2)\Rightarrow a_1=a_2
\]
for arbitrary $a_1,a_2\in A$.

\textbf{Available tools.}

Injectivity of $g$ says:
\[
g(b_1)=g(b_2)\Rightarrow b_1=b_2.
\]
Injectivity of $f$ says:
\[
f(a_1)=f(a_2)\Rightarrow a_1=a_2.
\]

\textbf{Intermediate states.}

Starting from
\[
g(f(a_1))=g(f(a_2)),
\]
injectivity of $g$ gives
\[
f(a_1)=f(a_2).
\]
Then injectivity of $f$ gives
\[
a_1=a_2.
\]

\textbf{Compressed proof.}

\begin{quote}
Suppose $(g\circ f)(a_1)=(g\circ f)(a_2)$. Then $g(f(a_1))=g(f(a_2))$, so $f(a_1)=f(a_2)$ by injectivity of $g$. By injectivity of $f$, $a_1=a_2$. Hence $g\circ f$ is injective.
\end{quote}

\subsection{What ``obvious'' often hides}

When a proof says ``clearly,'' it may be compressing one of several standard moves:

\begin{itemize}
\item applying a definition,
\item using a previously proved lemma,
\item doing algebra,
\item checking a domain condition,
\item handling a trivial case,
\item invoking a standard theorem.
\end{itemize}

Your job as a reader is to expand the word ``clearly'' into the actual move. Your job as a writer is to avoid compressing a step before you know exactly what it contains.

\begin{trapbox}
False intuition: short proofs are automatically better. A short proof is better only when the reader can reconstruct the omitted steps. Premature compression produces mystery, not elegance.
\end{trapbox}

\subsection{Proof sketch culture}

Advanced mathematics often uses proof sketches. A proof sketch is not a fake proof; it is a map of the main route, leaving routine checks to the reader. But what counts as routine depends on audience.

For a beginner, ``routine'' may include too much. For an expert, writing every algebraic rearrangement can hide the main idea. Mathematical writing is calibrated to a community.

\begin{labbox}
Compression drill. Take a one-page proof you have written. Produce three versions:
\begin{enumerate}
\item expanded beginner version, with every definition opened;
\item textbook version, with standard moves named;
\item proof sketch, with only the strategic spine.
\end{enumerate}
Then mark which audience each version serves.
\end{labbox}

\subsection{Exercises}

\begin{enumerate}
\item Prove that the product of two even integers is even.
\item Prove that the sum of an even integer and an odd integer is odd.
\item Prove that if $a\mid b$ and $a\mid c$, then $a\mid(b-c)$.
\item Prove that if $a\mid b$, then $a\mid bc$ for every integer $c$.
\item Prove that if $a,b,c\in\Z$, $a\mid b$, and $b\mid c$, then $a\mid c$.
\item Prove that the product of two rational numbers is rational.
\item Let $A,B,C$ be sets. Prove that if $A\subseteq B$ and $B\subseteq C$, then $A\subseteq C$.
\item Let $A$ and $B$ be sets. Prove that $A\cap B\subseteq A\cup B$.
\item Give a definition-machine proof that if a sequence is bounded above by $M$ and bounded below by $m$, then every term lies in the interval $[m,M]$.
\item Write a bad proof of one of the previous results, then annotate exactly where the definition machine was misused.
\end{enumerate}

\newpage
\section{If-Then Statements and Proof Strategy}

Most theorems are promises of the form
\[
P\Rightarrow Q.
\]
This is read ``if $P$, then $Q$.'' The statement $P$ is the hypothesis. The statement $Q$ is the conclusion.

The proof task is not to prove $P$. It is not to prove $Q$ in empty space. The task is to prove that $Q$ must hold in every situation where $P$ holds.

\begin{focusbox}
An implication is a bridge. You do not prove that everyone is standing at the entrance. You prove that anyone who starts at the entrance can get to the exit.
\end{focusbox}

This distinction is subtle at first, but it explains many proof strategies. A direct proof starts at $P$ and walks to $Q$. A contrapositive proof starts at $\neg Q$ and walks to $\neg P$. A contradiction proof assumes the theorem fails and shows the assumption breaks the system.

\subsection{The anatomy of an implication}

Consider:
\[
\text{If }n\text{ is divisible by }12,\text{ then }n\text{ is divisible by }3.
\]

The hypothesis is:
\[
n\text{ is divisible by }12.
\]

The conclusion is:
\[
n\text{ is divisible by }3.
\]

A proof should begin by assuming the hypothesis:

\begin{quote}
Suppose $n$ is divisible by $12$.
\end{quote}

Then it should aim at the conclusion:

\begin{quote}
We will show $n$ is divisible by $3$.
\end{quote}

Now the definitions tell us what to do. ``Divisible by $12$'' means
\[
n=12k
\]
for some integer $k$. ``Divisible by $3$'' means we want
\[
n=3m
\]
for some integer $m$.

Since $12k=3(4k)$, the witness is $m=4k$.

\begin{proof}
Suppose $n$ is divisible by $12$. Then $n=12k$ for some $k\in\Z$. Since
\[
n=12k=3(4k)
\]
and $4k\in\Z$, it follows that $n$ is divisible by $3$.
\end{proof}

That is direct proof in its cleanest form: assume the input condition, unfold it, and build the output condition.

\subsection{Truth of an implication}

An implication
\[
P\Rightarrow Q
\]
is false only in one case: $P$ is true and $Q$ is false.

This matches ordinary promises. Suppose a repair shop says:

\begin{quote}
If you bring the car before noon, we will inspect it today.
\end{quote}

The promise is broken only if you bring the car before noon and they do not inspect it today. If you do not bring the car before noon, the promise was not tested.

Mathematical implications work the same way. The statement
\[
\text{If }n\text{ is divisible by }12,\text{ then }n\text{ is divisible by }3
\]
does not claim that every integer is divisible by $12$. It says that every integer inside the hypothesis zone is also inside the conclusion zone.

\begin{formalbox}
The implication $P\Rightarrow Q$ is equivalent to saying:
\[
\text{there is no case where }P\text{ is true and }Q\text{ is false.}
\]
\end{formalbox}

This is why counterexamples to implications have a specific shape. To disprove
\[
P\Rightarrow Q,
\]
you must find an example where $P$ holds and $Q$ fails.

\subsection{Vacuous truth without the fog}

The statement
\[
\text{If }x\in\emptyset,\text{ then }x^2=17
\]
is true for every real number $x$ if we read it as a universal implication:
\[
\forall x\in\R,\ x\in\emptyset\Rightarrow x^2=17.
\]
Why? Because there is no $x\in\emptyset$. There is no case with a true hypothesis and false conclusion.

This can feel like a trick, but it is not a trick. It is bookkeeping. An implication says that all cases satisfying the hypothesis also satisfy the conclusion. If there are no cases satisfying the hypothesis, the containment is empty and therefore unbroken.

A physical analogy: suppose a sign says,

\begin{quote}
All passengers in row 99 must wear blue wristbands.
\end{quote}

If the theater has no row 99, the rule has not been violated. There is nobody in row 99 without a wristband.

\begin{trapbox}
Vacuous truth does not mean the conclusion is magically true. It means the conditional promise has no counterexample because the hypothesis never occurs.
\end{trapbox}

\subsection{Direct proof as forward engineering}

Direct proof is the first strategy to try when the hypothesis contains useful structure and the conclusion asks for structure you can build from it.

The pattern:

\begin{quote}
Assume $P$. Use definitions, known facts, and algebra to prove $Q$.
\end{quote}

Here is a classic example.

\begin{theorem}
If $n$ is an odd integer, then $n^2$ is odd.
\end{theorem}

\begin{proof}
Suppose $n$ is odd. Then there exists $k\in\Z$ such that $n=2k+1$. Squaring,
\[
n^2=(2k+1)^2=4k^2+4k+1=2(2k^2+2k)+1.
\]
Since $2k^2+2k\in\Z$, the number $n^2$ is odd.
\end{proof}

The hypothesis ``$n$ is odd'' opens into a useful expression. The conclusion ``$n^2$ is odd'' asks for an expression of the form $2(\text{integer})+1$. Direct proof is natural.

\begin{intuitionbox}
Direct proof feels like starting with raw material and machining it into the required shape.
\end{intuitionbox}

\subsection{When direct proof gets stuck}

Consider:

\begin{theorem}
If $n^2$ is even, then $n$ is even.
\end{theorem}

A direct proof would start:

\begin{quote}
Suppose $n^2$ is even. Then $n^2=2k$ for some integer $k$.
\end{quote}

We want to prove $n=2m$ for some integer $m$.

Can we take square roots and write
\[
n=\sqrt{2k}?
\]
That does not prove $n$ is twice an integer. The expression $\sqrt{2k}$ is usually not helpful.

The hypothesis gives information about $n^2$, but the conclusion is about $n$. Direct proof is possible with a stronger theorem about prime divisibility, but at this stage there is a cleaner strategy: prove the contrapositive.

\subsection{Contrapositive}

The contrapositive of
\[
P\Rightarrow Q
\]
is
\[
\neg Q\Rightarrow\neg P.
\]

These two statements are logically equivalent. They always have the same truth value.

For the theorem
\[
\text{If }n^2\text{ is even, then }n\text{ is even},
\]
the contrapositive is:
\[
\text{If }n\text{ is not even, then }n^2\text{ is not even}.
\]

For integers, ``not even'' means odd. So we prove:
\[
\text{If }n\text{ is odd, then }n^2\text{ is odd}.
\]

We already know how to prove that directly.

\begin{proof}
We prove the contrapositive. Suppose $n$ is not even. Since $n$ is an integer, $n$ is odd. Then $n=2k+1$ for some $k\in\Z$. Hence
\[
n^2=(2k+1)^2=4k^2+4k+1=2(2k^2+2k)+1.
\]
Thus $n^2$ is odd, so $n^2$ is not even. Therefore, by contrapositive, if $n^2$ is even, then $n$ is even.
\end{proof}

The contrapositive turns the awkward direction into the natural one.

\begin{formalbox}
\[
P\Rightarrow Q
\]
and
\[
\neg Q\Rightarrow\neg P
\]
are equivalent. The converse
\[
Q\Rightarrow P
\]
is a different statement. The inverse
\[
\neg P\Rightarrow\neg Q
\]
is also a different statement.
\end{formalbox}

\subsection{Converse confusion}

Beginners often prove the converse by accident.

Original theorem:
\[
\text{If }n^2\text{ is even, then }n\text{ is even}.
\]

Converse:
\[
\text{If }n\text{ is even, then }n^2\text{ is even}.
\]

The converse happens to be true here, but proving it does not prove the original. In many cases the converse is false.

Original:
\[
\text{If }n\text{ is divisible by }12,\text{ then }n\text{ is divisible by }3.
\]

Converse:
\[
\text{If }n\text{ is divisible by }3,\text{ then }n\text{ is divisible by }12.
\]

False. Take $n=6$. Then $6$ is divisible by $3$, but not by $12$.

\begin{trapbox}
When proving an implication, keep the direction visible. Write the hypothesis and conclusion separately before you begin.
\end{trapbox}

\subsection{Software analogy: preconditions and postconditions}

An implication behaves like a function contract.

\begin{quote}
Precondition: the input array is sorted.

Postcondition: binary search returns the index of the target if it exists.
\end{quote}

The contract does not say every array is sorted. It says: under this precondition, the algorithm guarantees this postcondition.

If you test binary search on an unsorted array and get nonsense, that does not refute the contract. You violated the hypothesis. If you test it on a sorted array and it fails to find an existing target, then you have a real counterexample.

The same discipline applies to theorems. When a theorem says
\[
P\Rightarrow Q,
\]
the hypothesis $P$ is the precondition. The conclusion $Q$ is the postcondition.

\subsection{Strategy selection}

Before writing a proof, pause and classify the statement.

\begin{itemize}
\item Is it a universal statement? You probably start with an arbitrary object.
\item Is it an existential statement? You probably need to build or identify a witness.
\item Is it an implication? You probably assume the hypothesis.
\item Is the conclusion a definition? You probably need to close that definition.
\item Is the conclusion negative? Contradiction or contrapositive may help.
\item Does the hypothesis look weak but the negated conclusion looks concrete? Try contrapositive.
\end{itemize}

Proof strategy is not a bag of tricks. It is reading the grammar of the statement and letting the grammar tell you what the first move should be.

\begin{labbox}
For each theorem below, write only the first proof sentence and the final target shape.

\begin{enumerate}
\item If $n$ is even, then $n^2+4n+6$ is even.
\item If $a\mid b$ and $b\mid c$, then $a\mid c$.
\item If $A\subseteq B$ and $B\subseteq C$, then $A\subseteq C$.
\item If $x$ is rational and $y$ is rational, then $x-y$ is rational.
\end{enumerate}
\end{labbox}

\subsection{Worked proof: rational minus rational}

\begin{theorem}
If $x$ and $y$ are rational numbers, then $x-y$ is rational.
\end{theorem}

\begin{proof}
Suppose $x$ and $y$ are rational. Then there exist integers $a,b,c,d$ with $b\neq0$ and $d\neq0$ such that
\[
x=\frac ab,\qquad y=\frac cd.
\]
Then
\[
x-y=\frac ab-\frac cd=\frac{ad-bc}{bd}.
\]
Since $ad-bc\in\Z$, $bd\in\Z$, and $bd\neq0$, the number $x-y$ is rational.
\end{proof}

The hypothesis gives two rational witnesses. The conclusion asks for one rational witness for $x-y$. The proof builds it:
\[
\frac{ad-bc}{bd}.
\]

\subsection{Worked proof: subset transitivity}

\begin{theorem}
If $A\subseteq B$ and $B\subseteq C$, then $A\subseteq C$.
\end{theorem}

\begin{proof}
Suppose $A\subseteq B$ and $B\subseteq C$. Let $x\in A$ be arbitrary. Since $A\subseteq B$ and $x\in A$, we have $x\in B$. Since $B\subseteq C$ and $x\in B$, we have $x\in C$. Therefore every element of $A$ is an element of $C$, so $A\subseteq C$.
\end{proof}

The proof is short because subset inclusion is a transport machine:
\[
x\in A\quad\Rightarrow\quad x\in B\quad\Rightarrow\quad x\in C.
\]

\subsection{Proof by cases}

Sometimes an implication has a hypothesis that naturally splits into alternatives. Then proof by cases is appropriate.

\begin{theorem}
For every integer $n$, the product $n(n+1)$ is even.
\end{theorem}

\begin{proof}
Let $n\in\Z$ be arbitrary. Every integer is either even or odd.

Case 1: $n$ is even. Then $n=2k$ for some $k\in\Z$, so
\[
n(n+1)=2k(n+1)=2(k(n+1)).
\]
Since $k(n+1)\in\Z$, the product is even.

Case 2: $n$ is odd. Then $n=2k+1$ for some $k\in\Z$. Therefore $n+1=2k+2=2(k+1)$, so
\[
n(n+1)=n\cdot2(k+1)=2(n(k+1)).
\]
Since $n(k+1)\in\Z$, the product is even.

In both cases, $n(n+1)$ is even.
\end{proof}

The cases are exhaustive because every integer is even or odd. They are useful because each case provides a different form that makes the conclusion easy.

\begin{trapbox}
In proof by cases, the cases must cover every possibility allowed by the hypothesis. A proof that checks only convenient cases is just a partial proof.
\end{trapbox}

\subsection{Biconditionals}

A biconditional
\[
P\Longleftrightarrow Q
\]
means both:
\[
P\Rightarrow Q
\]
and
\[
Q\Rightarrow P.
\]

The usual proof structure is therefore two implications.

\begin{theorem}
For an integer $n$, $n$ is even if and only if $n^2$ is even.
\end{theorem}

\begin{proof}
First suppose $n$ is even. Then $n=2k$ for some $k\in\Z$. Hence
\[
n^2=(2k)^2=4k^2=2(2k^2).
\]
Since $2k^2\in\Z$, $n^2$ is even.

Conversely, suppose $n^2$ is even. We prove this direction by contrapositive. If $n$ is not even, then $n$ is odd. So $n=2k+1$ for some $k\in\Z$. Then
\[
n^2=(2k+1)^2=2(2k^2+2k)+1,
\]
which is odd. Thus $n^2$ is not even. Therefore if $n^2$ is even, then $n$ is even.

Combining the two directions, $n$ is even if and only if $n^2$ is even.
\end{proof}

The phrase ``if and only if'' is strong. It says the two conditions are perfectly synchronized.

\begin{intuitionbox}
A biconditional says two tests recognize exactly the same objects.
\end{intuitionbox}

\subsection{Bad proof debugged: proving the wrong direction}

Claim:

\begin{theorem}
If $n^2$ is odd, then $n$ is odd.
\end{theorem}

Bad proof:

\begin{quote}
Suppose $n$ is odd. Then $n=2k+1$. Squaring gives $n^2=2(2k^2+2k)+1$, so $n^2$ is odd. Therefore if $n^2$ is odd, then $n$ is odd.
\end{quote}

The proof establishes:
\[
n\text{ odd}\Rightarrow n^2\text{ odd}.
\]
That is the converse of the desired statement.

A correct proof can use the contrapositive:

\begin{proof}
We prove the contrapositive. Suppose $n$ is not odd. Since $n$ is an integer, $n$ is even. Then $n=2k$ for some $k\in\Z$. Hence
\[
n^2=4k^2=2(2k^2).
\]
So $n^2$ is even, and therefore $n^2$ is not odd. By contrapositive, if $n^2$ is odd, then $n$ is odd.
\end{proof}

The repair is not longer because mathematics enjoys ceremony. It is longer because direction matters.

\subsection{History and connection: implication as structure}

\begin{historybox}
The formal study of implication goes back to ancient logic, especially Aristotle's work on syllogisms, but the modern symbolic treatment grew through Boole, Frege, Peano, Russell, Hilbert, and others. In the twentieth century, implication became central not only in mathematics but also in computing: proof systems, type systems, program verification, and automated theorem provers all depend on precise conditional structure.
\end{historybox}

A theorem is a reusable implication. Once proved, it becomes a component. If a later proof can establish the theorem's hypothesis, it can import the theorem's conclusion.

This is exactly how large mathematical theories scale. You do not reprove everything. You prove reliable implications and compose them.

\subsection{A proof-strategy checklist}

When facing a new theorem, write:

\[
\text{Hypotheses:}\quad \cdots
\]
\[
\text{Conclusion:}\quad \cdots
\]

Then ask:

\begin{enumerate}
\item What definitions can I open from the hypotheses?
\item What definition must I close for the conclusion?
\item Is the statement universal? If yes, what arbitrary object should I choose?
\item Is the statement existential? If yes, what witness should I build?
\item Is direct proof natural?
\item Would the contrapositive be clearer?
\item Would contradiction turn the target into a useful false assumption?
\item Are there cases that cover all possibilities?
\end{enumerate}

This checklist is not training wheels. Experienced mathematicians do a compressed version of it mentally.

\subsection{Exercises}

\begin{enumerate}
\item Prove directly: if $n$ is even, then $5n+8$ is even.
\item Prove directly: if $n$ is odd, then $3n+2$ is odd.
\item Prove by contrapositive: if $n^2$ is odd, then $n$ is odd.
\item Prove by cases: for every integer $n$, $n^2+n$ is even.
\item Prove: if $a\mid b$ and $c\mid d$, then $ac\mid bd$.
\item Prove: if $x$ and $y$ are rational, then $3x-5y$ is rational.
\item Prove: if $A\subseteq B$, then $A\cap C\subseteq B\cap C$.
\item Find a counterexample to the converse: if $n$ is divisible by $4$, then $n$ is even.
\item For each statement in this section, identify whether the proof used direct proof, contrapositive, cases, or a combination.
\item Write a one-paragraph strategy for proving: if $n^2-1$ is not divisible by $8$, then $n$ is even.
\end{enumerate}

\newpage
\section{Proof by Contradiction}

Proof by contradiction is one of the most dramatic methods in mathematics. It begins by denying what we want to prove. Then it follows that denial carefully until something impossible appears.

The basic form is:

\[
\text{To prove }P,\text{ assume }\neg P\text{ and derive a contradiction.}
\]

Once the contradiction appears, the assumption $\neg P$ cannot be maintained. Therefore $P$ must be true.

\begin{focusbox}
Contradiction is not a license to be vague. The contradiction must be a precise collision: a statement and its negation, or an impossible fact such as $0=1$, an integer being both even and odd, or a positive number being less than itself.
\end{focusbox}

\subsection{A tiny contradiction proof}

\begin{theorem}
There is no smallest positive real number.
\end{theorem}

\begin{proof}
Suppose, for contradiction, that there is a smallest positive real number $r$. Since $r>0$, the number $r/2$ is also positive. But
\[
0<\frac r2<r.
\]
This contradicts the assumption that $r$ is the smallest positive real number. Therefore there is no smallest positive real number.
\end{proof}

The impossible collision is exact:

\begin{itemize}
\item Assumption: $r$ is the smallest positive real number.
\item Derived fact: $r/2$ is a positive real number smaller than $r$.
\end{itemize}

Those cannot both be true.

\subsection{Contradiction versus contrapositive}

Contradiction and contrapositive are related, but they feel different in practice.

To prove an implication
\[
P\Rightarrow Q
\]
by contrapositive, prove:
\[
\neg Q\Rightarrow\neg P.
\]

To prove it by contradiction, assume:
\[
P\text{ and }\neg Q
\]
and derive an impossibility.

For example, to prove
\[
\text{if }n^2\text{ is even, then }n\text{ is even},
\]
contrapositive assumes $n$ is odd and proves $n^2$ is odd.

Contradiction assumes $n^2$ is even and $n$ is odd, then derives that $n^2$ is odd, contradicting that $n^2$ is even.

Both are valid. The contrapositive proof often feels cleaner because it produces a positive conclusion rather than carrying a contradiction flag.

\begin{intuitionbox}
Contrapositive says, ``If the destination is absent, the starting condition was absent.'' Contradiction says, ``Suppose the starting condition holds but the destination fails; now watch the system break.''
\end{intuitionbox}

\subsection{The contradiction contract}

A contradiction proof has a contract:

\begin{enumerate}
\item State what you are assuming for contradiction.
\item Use that assumption together with known facts.
\item Derive a clear impossibility.
\item Identify the impossibility.
\item Conclude that the contradictory assumption was false.
\end{enumerate}

Skipping step 4 is a common beginner mistake. A proof that ends with ``which is a contradiction'' but does not say what contradicted what leaves the reader to do the proof's final job.

\begin{trapbox}
Do not use contradiction as a fog machine. If you cannot name the contradictory pair, the proof is not finished.
\end{trapbox}

\subsection{Classic proof: irrationality of square root of 2}

\begin{theorem}
$\sqrt2$ is irrational.
\end{theorem}

\begin{proof}
Suppose, for contradiction, that $\sqrt2$ is rational. Then there exist integers $a,b$ with $b\neq0$ such that
\[
\sqrt2=\frac ab.
\]
Choose such a fraction in lowest terms, so $a$ and $b$ have no common factor greater than $1$.

Squaring both sides gives
\[
2=\frac{a^2}{b^2},
\]
so
\[
a^2=2b^2.
\]
Thus $a^2$ is even. Since $a^2$ is even, $a$ is even. Therefore $a=2k$ for some $k\in\Z$.

Substitute $a=2k$ into $a^2=2b^2$:
\[
(2k)^2=2b^2.
\]
So
\[
4k^2=2b^2,
\]
and hence
\[
b^2=2k^2.
\]
Thus $b^2$ is even, so $b$ is even.

We have shown that both $a$ and $b$ are even. Therefore they share a factor of $2$, contradicting the choice of $a/b$ in lowest terms. Hence $\sqrt2$ is irrational.
\end{proof}

This proof is famous because it turns a geometric fact into arithmetic pressure. The diagonal of a unit square has length $\sqrt2$. If that length were a ratio of whole numbers, the ratio could be chosen in lowest terms. But the equation forces both numerator and denominator to be even, so the ratio was not in lowest terms after all.

\begin{historybox}
The discovery of irrational magnitudes is traditionally linked to the Pythagorean school. The historical details are hard to separate from legend, but the mathematical shock is real: geometry produced lengths that ratios of whole numbers could not express. This forced a major expansion in the concept of number.
\end{historybox}

\subsection{Why lowest terms matters}

In the proof of irrationality of $\sqrt2$, the phrase ``choose the fraction in lowest terms'' is essential.

If we only assume
\[
\sqrt2=\frac ab
\]
for some integers $a,b$, then showing $a$ and $b$ are both even is not immediately impossible. Fractions are often not in lowest terms:
\[
\frac46=\frac23.
\]
The contradiction requires the extra condition that the representation has no common factor.

Why are we allowed to choose a lowest-terms representation? Because every rational number can be written as a fraction of integers with no common factor greater than $1$. That fact comes from the well-ordering principle or the Euclidean algorithm.

The contradiction is:

\[
\text{$a$ and $b$ have no common factor greater than $1$}
\]
and
\[
\text{$a$ and $b$ are both divisible by $2$}.
\]

\begin{formalbox}
A contradiction proof often needs a carefully strengthened assumption. ``Let $a/b$ be in lowest terms'' is not decoration; it creates the exact condition that will later be contradicted.
\end{formalbox}

\subsection{Bad proof debugged: square root of 2}

Bad proof:

\begin{quote}
Suppose $\sqrt2=a/b$. Then $a^2=2b^2$. So $a^2$ is even, hence $a$ is even. Similarly $b$ is even. This is impossible. Therefore $\sqrt2$ is irrational.
\end{quote}

There are two problems.

First, ``similarly $b$ is even'' is not immediate. You need to substitute $a=2k$ and derive $b^2=2k^2$.

Second, ``This is impossible'' is false unless the fraction was chosen in lowest terms. It is not impossible for two integers to both be even. It is impossible for numerator and denominator of a lowest-terms fraction to both be even.

Corrected proof skeleton:

\begin{enumerate}
\item Assume $\sqrt2=a/b$ in lowest terms.
\item Derive $a^2=2b^2$.
\item Conclude $a$ is even, so $a=2k$.
\item Substitute and conclude $b$ is even.
\item Contradict lowest terms.
\end{enumerate}

\subsection{Contradiction with divisibility}

\begin{theorem}
There is no integer $n$ such that $n$ is both even and odd.
\end{theorem}

\begin{proof}
Suppose, for contradiction, that there exists an integer $n$ that is both even and odd. Since $n$ is even, $n=2a$ for some $a\in\Z$. Since $n$ is odd, $n=2b+1$ for some $b\in\Z$. Thus
\[
2a=2b+1.
\]
Rearranging gives
\[
2(a-b)=1.
\]
This says that $1$ is even, since $a-b\in\Z$. But $1$ is not even. Contradiction. Therefore no integer is both even and odd.
\end{proof}

The contradiction here is not emotional. It is exact:
\[
1=2(a-b)
\]
would make $1$ divisible by $2$, which it is not.

\subsection{Existence proofs by contradiction}

Contradiction can prove existence. To prove that some object exists, assume no such object exists and show impossibility.

\begin{theorem}
There exist irrational numbers $x$ and $y$ such that $x+y$ is rational.
\end{theorem}

This theorem actually has a direct proof:
\[
x=\sqrt2,\qquad y=-\sqrt2.
\]
Both are irrational, but
\[
x+y=0,
\]
which is rational.

Contradiction is not needed.

This matters because beginners sometimes use contradiction whenever a theorem feels hard. But contradiction is not always the cleanest method. For existence claims, a witness is usually better than a contradiction if you can find one.

\begin{focusbox}
If you can prove existence by showing the object, show the object. A witness beats a thunderstorm.
\end{focusbox}

\subsection{Nonconstructive contradiction}

Some existence proofs by contradiction do not produce a concrete object. They show that nonexistence is impossible.

Here is a simple finite example.

\begin{theorem}
Among any three integers, two have the same parity.
\end{theorem}

\begin{proof}
Suppose, for contradiction, that among three integers no two have the same parity. Then at most one of the integers is even and at most one is odd. That accounts for at most two integers total. But we have three integers. Contradiction. Therefore at least two of the three integers have the same parity.
\end{proof}

This is a pigeonhole argument in miniature. There are two parity boxes, even and odd. Three integers cannot occupy two boxes without some box receiving at least two integers.

\begin{historybox}
The pigeonhole principle is often associated with Dirichlet, who used counting arguments to prove deep results in number theory. Its simplest form says that if more objects than boxes are distributed among the boxes, at least one box contains more than one object.
\end{historybox}

\subsection{Infinite primes}

Euclid's proof that there are infinitely many primes is one of the great contradiction proofs.

\begin{theorem}
There are infinitely many prime numbers.
\end{theorem}

\begin{proof}
Suppose, for contradiction, that there are only finitely many primes. List them as
\[
p_1,p_2,\ldots,p_n.
\]
Now form the number
\[
N=p_1p_2\cdots p_n+1.
\]
The integer $N$ is greater than $1$, so it has a prime divisor $q$.

Since $p_1,\ldots,p_n$ are assumed to be all the primes, $q$ must equal one of them, say $q=p_i$. Because $q$ divides the product $p_1p_2\cdots p_n$, and $q$ also divides $N$, it follows that $q$ divides the difference
\[
N-p_1p_2\cdots p_n=1.
\]
But no prime divides $1$. This is a contradiction. Therefore there are infinitely many primes.
\end{proof}

The engine is clever:

\[
N=p_1p_2\cdots p_n+1.
\]

This number is designed to leave remainder $1$ when divided by every prime on the supposed complete list. So it either is prime itself or has a prime factor not on the list. Either way, the list was incomplete.

\begin{intuitionbox}
Euclid's construction is like checking a supposedly complete registry, then building a new object that evades every entry in the registry.
\end{intuitionbox}

\subsection{The exact contradiction in Euclid's proof}

The contradiction is sometimes misstated as:

\begin{quote}
$N$ is not divisible by any prime, which is impossible.
\end{quote}

That is close, but the clean version is:

\begin{itemize}
\item Every integer greater than $1$ has a prime divisor.
\item If the list contains all primes, that divisor must be one of $p_1,\ldots,p_n$.
\item But no $p_i$ divides $N$, because $N$ leaves remainder $1$ when divided by $p_i$.
\end{itemize}

The contradiction is between ``$q$ is on the complete list'' and ``$q$ cannot be any prime on the list.''

\begin{trapbox}
Euclid's proof does not say $p_1p_2\cdots p_n+1$ is always prime. For $2\cdot3\cdot5\cdot7\cdot11\cdot13+1=30031$, the result is composite. The point is that its prime factors are not on the original list.
\end{trapbox}

\subsection{Contradiction and minimal counterexamples}

A powerful method combines contradiction with choosing a smallest counterexample.

The rough pattern:

\begin{enumerate}
\item Suppose the theorem is false.
\item Then there is at least one counterexample.
\item Choose the smallest counterexample.
\item Use it to build an even smaller counterexample.
\item Contradiction.
\end{enumerate}

Here is a gentle example.

\begin{theorem}
Every integer $n\geq2$ has a prime divisor.
\end{theorem}

\begin{proof}
Suppose, for contradiction, that some integer $n\geq2$ has no prime divisor. By the well-ordering principle, choose the smallest such integer and call it $m$.

If $m$ is prime, then $m$ has a prime divisor, namely itself. That contradicts the choice of $m$. Therefore $m$ is not prime.

Since $m\geq2$ and $m$ is not prime, it is composite. So $m=ab$ for integers $a,b$ with
\[
1<a<m,\qquad 1<b<m.
\]
Because $a<m$ and $m$ was the smallest integer at least $2$ with no prime divisor, the integer $a$ must have a prime divisor $p$. Since $p\mid a$ and $a\mid m$, we get $p\mid m$. This contradicts the choice of $m$ as having no prime divisor.

Therefore every integer $n\geq2$ has a prime divisor.
\end{proof}

Minimal counterexample arguments are common in number theory, combinatorics, graph theory, and induction. They are a formal way of saying: if failure happens, let us inspect the first failure. Often the first failure cannot exist.

\subsection{Danger: proving too much}

A contradiction proof can go wrong by using an assumption that is stronger than the negation of the theorem.

Suppose you want to prove:

\[
\text{There exists an integer }n\text{ such that }n^2>10.
\]

The negation is:
\[
\text{For every integer }n,\ n^2\leq10.
\]

It is not:
\[
\text{For every integer }n,\ n^2<0.
\]

If you assume something stronger than the actual negation and derive a contradiction, you have not proved the theorem. You only proved that your stronger assumption was false.

\begin{formalbox}
To prove $P$ by contradiction, assume exactly $\neg P$, or at least an assumption that follows from $\neg P$. Do not invent a stronger denial because it is easier to break.
\end{formalbox}

\subsection{Negating statements before contradiction}

Contradiction requires accurate negation.

Original:
\[
\forall x\in A,\ P(x).
\]
Negation:
\[
\exists x\in A\text{ such that }\neg P(x).
\]

Original:
\[
\exists x\in A\text{ such that }P(x).
\]
Negation:
\[
\forall x\in A,\ \neg P(x).
\]

Original:
\[
P\Rightarrow Q.
\]
Negation:
\[
P\text{ and }\neg Q.
\]

If the negation is wrong, the contradiction proof is proving the wrong thing.

\begin{labbox}
Write the negation you would assume for contradiction:

\begin{enumerate}
\item Every bounded sequence converges.
\item There exists a largest integer.
\item If $n^2$ is divisible by $3$, then $n$ is divisible by $3$.
\item For every $\eps>0$, there exists $N\in\N$ such that $n\geq N$ implies $|a_n-L|<\eps$.
\end{enumerate}
\end{labbox}

\subsection{Bad proof debugged: no largest integer}

Claim:

\begin{theorem}
There is no largest integer.
\end{theorem}

Bad proof:

\begin{quote}
Suppose there is a largest integer $N$. Then $N+1$ is larger. Contradiction.
\end{quote}

This is almost fine, but it skips one small piece: $N+1$ must be an integer. The largest-integer assumption concerns integers, not all real numbers. A polished proof says:

\begin{proof}
Suppose, for contradiction, that there is a largest integer $N$. Since $N\in\Z$, we have $N+1\in\Z$. But $N+1>N$, contradicting the assumption that $N$ is the largest integer. Therefore there is no largest integer.
\end{proof}

This may feel picky. It is picky in the same way bridge inspection is picky. The missing sentence is exactly what connects the constructed object back to the domain of the claim.

\subsection{Contradiction in later mathematics}

Contradiction appears everywhere.

In real analysis, it proves uniqueness of limits:

\begin{quote}
Assume a sequence has two different limits $L$ and $M$. Choose disjoint neighborhoods around $L$ and $M$. Eventually the sequence must be inside both, impossible.
\end{quote}

In linear algebra, it proves uniqueness of inverses:

\begin{quote}
Assume a matrix has two inverses $B$ and $C$. Then $B=BI=B(AC)=(BA)C=IC=C$.
\end{quote}

In topology, it proves separation results. In algebra, it proves nonexistence results. In computability, diagonal arguments often have contradiction at their core.

\begin{historybox}
Proof by contradiction is classically accepted in standard mathematics, but it has also been philosophically debated. Constructivist mathematicians distinguish between proving that nonexistence is impossible and actually constructing an object. This debate is not just philosophy; it affects proof assistants, type theory, and the meaning of existence in computational mathematics.
\end{historybox}

\subsection{Walking questions}

\begin{walkbox}
When you see a contradiction proof, ask:

\begin{enumerate}
\item What exact statement is being denied?
\item Is the denial correct?
\item Which assumptions are used after the denial?
\item What is the contradictory pair?
\item Could the proof be rewritten as a contrapositive or direct proof?
\end{enumerate}
\end{walkbox}

\subsection{Exercises}

\begin{enumerate}
\item Prove by contradiction that there is no largest real number.
\item Prove by contradiction that there is no smallest positive integer greater than $1$ that is divisible by every positive integer.
\item Prove that $\sqrt3$ is irrational. You may use the fact that if $3\mid n^2$, then $3\mid n$.
\item Prove that $\sqrt5$ is irrational. You may use the analogous divisibility fact.
\item Prove that there do not exist $a,b\in\Z$ such that $a^2-4b=2$.
\item Prove there are infinitely many integers divisible by $7$.
\item Prove that if $x+y=0$ and $x>0$, then $y<0$ by contradiction.
\item Write a contradiction proof of: if $A\subseteq B$ and $x\notin B$, then $x\notin A$.
\item Find the flaw in this proof: ``Suppose $n$ is the smallest integer. Then $n-1<n$, contradiction. Therefore no integer is smallest.'' Repair it.
\item For each contradiction proof in this section, identify the exact contradiction in one sentence.
\end{enumerate}

\newpage
\section{Examples, Non-examples, and Counterexamples}

Examples are not the snack after the theorem. They are part of the meal.

A definition without examples is a machine you have never seen run. A theorem without examples is a road sign for a road you have never walked. Counterexamples are stress tests. Non-examples are boundary markers. Together they build mathematical taste.

\begin{focusbox}
Examples teach you what a definition includes. Non-examples teach you what it excludes. Counterexamples teach you how statements fail.
\end{focusbox}

This section is about a habit: whenever you meet a new concept, build a small zoo around it. Do not collect only pretty examples. Collect edge cases, boring cases, tiny cases, broken cases, and almost-cases.

\subsection{Examples versus proofs}

An example can prove an existential statement.

Statement:
\[
\exists n\in\Z\text{ such that }n^2=25.
\]

Proof:
\[
n=5.
\]

That is enough. An existential statement asks for at least one witness. A single valid witness proves it.

But an example cannot prove a universal statement.

Statement:
\[
\forall n\in\Z,\ n^2\geq0.
\]

Checking $n=5$ is not a proof. Checking $n=-100$ is not a proof. Checking a million integers by computer is evidence, but not proof. The statement concerns every integer, and there are infinitely many.

To prove it, use a general argument:

\begin{proof}
Let $n\in\Z$ be arbitrary. If $n\geq0$, then $n^2\geq0$. If $n<0$, then $-n>0$, and $n^2=(-n)^2>0$. Therefore $n^2\geq0$ for every integer $n$.
\end{proof}

The proof handles all integers at once because it reasons from arbitrary structure, not from sampled values.

\begin{trapbox}
An example proves ``there exists.'' A counterexample disproves ``for all.'' Neither one proves a universal statement.
\end{trapbox}

\subsection{Non-examples}

A non-example is an object that fails a definition.

Definition:
\[
n\text{ is even}\quad\Longleftrightarrow\quad n=2k\text{ for some }k\in\Z.
\]

Example: $10$ is even because $10=2\cdot5$.

Non-example: $11$ is not even because there is no integer $k$ such that $11=2k$.

Non-examples are useful because they force you to use the exact definition. Saying ``$11$ feels odd'' is not proof. Saying ``if $11=2k$, then $k=11/2$, which is not an integer'' is proof.

\begin{intuitionbox}
A non-example is not merely absent from the club. It fails the entry test, and you should know which part of the test it fails.
\end{intuitionbox}

\subsection{Counterexamples}

A counterexample is an example that disproves a claim.

Claim:
\[
\forall x,y\in\R,\ x^2=y^2\Rightarrow x=y.
\]

Counterexample:
\[
x=1,\qquad y=-1.
\]

Then
\[
x^2=1=y^2,
\]
but
\[
x\neq y.
\]

The counterexample has the exact shape needed to refute the implication:

\[
\text{hypothesis true, conclusion false.}
\]

For a universal implication
\[
\forall x,\ P(x)\Rightarrow Q(x),
\]
a counterexample must provide an $x$ such that $P(x)$ is true and $Q(x)$ is false.

\begin{formalbox}
To disprove
\[
\forall x\in D,\ P(x),
\]
find $x\in D$ such that $\neg P(x)$.

To disprove
\[
\forall x\in D,\ P(x)\Rightarrow Q(x),
\]
find $x\in D$ such that $P(x)$ is true and $Q(x)$ is false.
\end{formalbox}

\subsection{The domain is part of the counterexample}

Claim:
\[
\forall x,\ x^2\geq x.
\]

Is this true? It depends on the domain.

If $x\in\Z$, the statement is true. For integers, either $x\leq0$, in which case $x^2\geq0\geq x$, or $x\geq1$, in which case $x^2\geq x$.

If $x\in\R$, the statement is false. Take
\[
x=\frac12.
\]
Then
\[
x^2=\frac14<\frac12=x.
\]

So the counterexample is not just ``$1/2$.'' It is ``$1/2$ in the real-number domain.'' If the theorem was about integers, $1/2$ is not even eligible.

\begin{trapbox}
A counterexample must live inside the stated domain. An object outside the domain refutes nothing.
\end{trapbox}

\subsection{Near misses}

A near miss almost satisfies a definition but fails one condition. Near misses are valuable because they reveal why each condition exists.

Definition:

\begin{quote}
A function $f:A\to B$ is injective if for all $x,y\in A$, $f(x)=f(y)$ implies $x=y$.
\end{quote}

Example:
\[
f:\R\to\R,\qquad f(x)=2x+1
\]
is injective.

Non-example:
\[
g:\R\to\R,\qquad g(x)=x^2
\]
is not injective because $g(1)=g(-1)$ but $1\neq-1$.

Near miss:
\[
h:[0,\infty)\to\R,\qquad h(x)=x^2
\]
is injective. Same formula, different domain.

This near miss teaches a serious lesson: the function is not just the formula. The domain and codomain are part of the object.

\subsection{Injective and surjective examples}

Let
\[
f:\N\to\N,\qquad f(n)=n+1,
\]
where $\N=\{1,2,3,\ldots\}$.

This function is injective. If
\[
f(a)=f(b),
\]
then
\[
a+1=b+1,
\]
so $a=b$.

It is not surjective. There is no $n\in\N$ such that
\[
n+1=1.
\]
The codomain element $1$ is missed.

Now let
\[
g:\{1,2\}\to\{a\}
\]
be defined by
\[
g(1)=a,\qquad g(2)=a.
\]
This function is surjective because the only codomain element, $a$, is hit. It is not injective because two different inputs have the same output.

These examples separate two ideas that beginners often blur:

\begin{itemize}
\item Injective means no merging.
\item Surjective means no missing codomain elements.
\end{itemize}

\subsection{Counterexamples to converses}

If a theorem says
\[
P\Rightarrow Q,
\]
the converse is
\[
Q\Rightarrow P.
\]

The converse may be true or false. A counterexample to the converse must have $Q$ true and $P$ false.

Theorem:
\[
\text{If }n\text{ is divisible by }12,\text{ then }n\text{ is divisible by }3.
\]

Converse:
\[
\text{If }n\text{ is divisible by }3,\text{ then }n\text{ is divisible by }12.
\]

Counterexample:
\[
n=6.
\]

Then $6$ is divisible by $3$, but $6$ is not divisible by $12$.

This is a clean counterexample because it satisfies the converse hypothesis and fails the converse conclusion.

\begin{walkbox}
For any implication you meet today, say its converse out loud. Then ask whether the converse is true, false, or uncertain. If false, try to build a counterexample.
\end{walkbox}

\subsection{Counterexamples to overgeneralization}

A common source of false statements is taking a true theorem and removing a hypothesis.

True theorem:

\begin{quote}
If $f$ is continuous on a closed bounded interval $[a,b]$, then $f$ attains a maximum on $[a,b]$.
\end{quote}

Remove ``closed'':
\[
f(x)=x\quad\text{on }(0,1).
\]
The function is continuous and bounded, but it does not attain a maximum. Its supremum is $1$, but $1$ is not in the domain.

Remove ``bounded interval'':
\[
f(x)=x\quad\text{on }[0,\infty).
\]
The function is continuous on a closed set, but it has no maximum.

Remove ``continuous'':
\[
f(x)=
\begin{cases}
x, & 0\leq x<1,\\
0, & x=1
\end{cases}
\quad\text{on }[0,1].
\]
This function is defined on a closed bounded interval, but it does not attain its supremum $1$.

The theorem's assumptions are not fancy wrapping. They are load-bearing.

\begin{historybox}
Nineteenth-century analysis was shaped by examples that broke plausible statements about functions. Mathematicians learned that continuity, compactness, differentiability, and convergence needed precise hypotheses. Counterexamples did not weaken analysis; they made it stronger.
\end{historybox}

\subsection{Bad counterexamples}

Claim:
\[
\text{Every prime number is odd.}
\]

A student says:

\begin{quote}
Counterexample: $9$, because $9$ is odd but not prime.
\end{quote}

This is not a counterexample. The claim has the form:
\[
\text{prime}\Rightarrow\text{odd}.
\]
To refute it, we need:
\[
\text{prime and not odd}.
\]
The correct counterexample is $2$, because $2$ is prime and even.

Another bad counterexample:

Claim:
\[
\forall x\in\R,\ x^2+1>0.
\]

Counterexample offered:
\[
x=i.
\]

But $i\notin\R$. The proposed object is outside the domain. It does not refute the real-number claim.

\begin{trapbox}
For an implication, a counterexample must satisfy the hypothesis. For a quantified statement, a counterexample must belong to the quantified domain.
\end{trapbox}

\subsection{Examples that prove existence}

Statement:
\[
\exists a,b\in\Z\text{ such that }a^2-b^2=15.
\]

Try small numbers:
\[
4^2-1^2=16-1=15.
\]
So $a=4$, $b=1$ proves the statement.

Statement:
\[
\exists x\in\R\text{ such that }x^2=2.
\]

Take $x=\sqrt2$. This proves existence in $\R$ because $\sqrt2$ is a real number.

Statement:
\[
\exists x\in\Q\text{ such that }x^2=2.
\]

The same candidate does not work because $\sqrt2\notin\Q$. In fact the statement is false.

Again, the domain matters.

\subsection{Families of examples}

One example can prove existence, but a family of examples builds understanding.

Even integers:
\[
\ldots,-4,-2,0,2,4,\ldots
\]
Family form:
\[
2k,\qquad k\in\Z.
\]

Odd integers:
\[
\ldots,-3,-1,1,3,5,\ldots
\]
Family form:
\[
2k+1,\qquad k\in\Z.
\]

Multiples of $7$:
\[
7k,\qquad k\in\Z.
\]

Rational numbers:
\[
\frac pq,\qquad p,q\in\Z,\ q\neq0.
\]

Upper bounds for $(1-1/n)$:
\[
1,\ 2,\ 100,\ \pi
\]
are all upper bounds. The least upper bound is $1$, but being an upper bound is a weaker condition.

Families reveal the degrees of freedom in a definition. They show what can vary while the structure remains.

\subsection{Boundary examples}

Boundary examples live at the edge of a definition.

Is $0$ even? Yes:
\[
0=2\cdot0.
\]

Is $1$ prime? No. A prime number has exactly two positive divisors, $1$ and itself. The number $1$ has only one positive divisor.

Is the empty set a subset of every set? Yes. To prove $\emptyset\subseteq A$, we must show every element of $\emptyset$ belongs to $A$. There are no elements of $\emptyset$, so there is no counterexample.

Is the empty function from $\emptyset$ to $B$ a function? Yes. It assigns to every element of $\emptyset$ exactly one element of $B$, because there are no elements requiring assignment. This is another appearance of vacuous truth.

Boundary examples often feel strange because they test whether you are following the definition or relying on a picture.

\begin{intuitionbox}
If an example feels silly but satisfies the definition, keep it. Silly examples are excellent stress tests.
\end{intuitionbox}

\subsection{The role of small finite models}

When learning set theory, relations, functions, or algebra, small finite examples are priceless.

Let
\[
A=\{1,2,3\},\qquad B=\{2,3,4\}.
\]
Then
\[
A\cap B=\{2,3\},
\]
\[
A\cup B=\{1,2,3,4\},
\]
\[
A\setminus B=\{1\}.
\]

These examples are not deep, but they calibrate the definitions.

For a relation example, define $R$ on $\{1,2,3\}$ by
\[
xRy\quad\Longleftrightarrow\quad x\leq y.
\]
This relation is reflexive and transitive, but not symmetric.

Define $S$ by
\[
xSy\quad\Longleftrightarrow\quad x-y\text{ is even}.
\]
This relation is reflexive, symmetric, and transitive. It groups the set into parity classes.

Tiny models let you see properties without drowning in notation.

\subsection{Counterexample search strategies}

When a statement seems false, search systematically.

First, identify the shape of failure.

For
\[
\forall x,\ P(x)\Rightarrow Q(x),
\]
you need $P(x)$ true and $Q(x)$ false.

Second, try boundary cases:
\[
0,\quad1,\quad -1,\quad \emptyset,\quad \{0\},\quad \text{constant functions}.
\]

Third, try changing the domain or codomain. Many false statements about functions fail because of codomain choices.

Fourth, try degenerate examples:

\begin{itemize}
\item empty sets,
\item singleton sets,
\item zero maps,
\item identity maps,
\item constant sequences,
\item sequences with alternating signs.
\end{itemize}

Fifth, remove one assumption from a true theorem and see what breaks.

\begin{labbox}
Find counterexamples to these claims. For each one, write the hypothesis true/false and conclusion true/false.

\begin{enumerate}
\item If $ab=0$, then $a=0$ and $b=0$.
\item If $x^2>4$, then $x>2$.
\item If a function is not injective, then it is not surjective.
\item If $A\cup B=A\cup C$, then $B=C$.
\item Every bounded sequence is convergent.
\end{enumerate}
\end{labbox}

\subsection{Worked counterexamples}

Claim:
\[
ab=0\Rightarrow a=0\text{ and }b=0.
\]

Counterexample:
\[
a=0,\qquad b=5.
\]
Then $ab=0$, but $b\neq0$, so the conclusion ``$a=0$ and $b=0$'' is false.

The true theorem is:
\[
ab=0\Rightarrow a=0\text{ or }b=0.
\]

Claim:
\[
x^2>4\Rightarrow x>2.
\]

Counterexample:
\[
x=-3.
\]
Then $x^2=9>4$, but $x>2$ is false.

The true conclusion would be:
\[
x>2\text{ or }x<-2.
\]

Claim:
\[
A\cup B=A\cup C\Rightarrow B=C.
\]

Counterexample:
\[
A=\{1,2\},\qquad B=\{1\},\qquad C=\{2\}.
\]
Then
\[
A\cup B=\{1,2\}=A\cup C,
\]
but
\[
B\neq C.
\]

The set $A$ hides the difference between $B$ and $C$ because it already contains both elements.

\subsection{Examples in proof discovery}

Examples are not only for checking final answers. They help you discover proofs.

Suppose you want to prove:
\[
\text{The sum of the first }n\text{ odd positive integers is }n^2.
\]

Compute:
\[
1=1,
\]
\[
1+3=4,
\]
\[
1+3+5=9,
\]
\[
1+3+5+7=16.
\]

The examples suggest squares. More importantly, they suggest a visual proof: start with a $1\times1$ square, then add an L-shaped border of $3$ dots, then one of $5$, then one of $7$. Each odd number grows the square by one layer.

Examples can also suggest induction:
\[
1+3+\cdots+(2n-1)=n^2.
\]
Then the next term is
\[
2(n+1)-1=2n+1,
\]
and
\[
n^2+(2n+1)=(n+1)^2.
\]

The examples did not prove the theorem, but they revealed the mechanism.

\subsection{Examples in debugging definitions}

Suppose someone defines a prime number as:

\begin{quote}
A number divisible only by $1$ and itself.
\end{quote}

This is close, but incomplete. What kind of number? Positive integer? Integer greater than $1$?

If negative integers are allowed, is $-5$ prime? Its divisors include $-1,1,-5,5$. If $1$ is allowed, the definition accidentally makes $1$ look prime unless we require exactly two positive divisors or require the number to be greater than $1$.

A cleaner definition:

\begin{definition}
A positive integer $p$ is prime if $p>1$ and its only positive divisors are $1$ and $p$.
\end{definition}

Examples force definitions to become honest. They expose missing domain restrictions and ambiguous words.

\begin{historybox}
Mathematics often advances when examples pressure definitions. The modern definition of function had to grow broad enough to include strange functions. The definition of continuity had to survive pathological examples. The definition of manifold had to separate local Euclidean behavior from global shape. Good definitions are not guessed fully formed; they are tested.
\end{historybox}

\subsection{Physical analogy: crash testing}

Engineers do not test a design only under pleasant conditions. They test load, heat, vibration, fatigue, edge cases, and failure modes.

Counterexamples are mathematical crash tests. A theorem that survives only nice examples is not a theorem yet. A definition that behaves well only for familiar examples is not mature yet.

This is why the empty set, zero, identity functions, constant functions, and weird sequences appear so often. They are not annoyances. They are test fixtures.

\subsection{Software analogy: unit tests and property tests}

An example is like a unit test:
\[
f(2)=4.
\]
It checks one case.

A theorem is like a property test:
\[
\forall x,\ f(x)\geq0.
\]
It describes a behavior over a whole domain.

A counterexample is a failing test case with a stack trace. It tells you not merely that something failed, but where:

\begin{itemize}
\item domain issue,
\item missing hypothesis,
\item wrong direction,
\item too-strong conclusion,
\item hidden edge case.
\end{itemize}

Good mathematicians keep a mental test suite. When they see a proposed theorem about functions, they test constant functions, identity functions, squaring, absolute value, step functions, and empty-domain functions. When they see a theorem about sets, they test empty sets, equal sets, disjoint sets, and nested sets.

\subsection{Non-examples in analysis}

Definition:

\begin{quote}
A sequence $(a_n)$ converges to $L$ if for every $\eps>0$, there exists $N\in\N$ such that for all $n\geq N$, $|a_n-L|<\eps$.
\end{quote}

Example:
\[
a_n=\frac1n
\]
converges to $0$.

Non-example:
\[
a_n=(-1)^n
\]
does not converge. It keeps jumping between $1$ and $-1$.

Near miss:
\[
a_n=\frac{(-1)^n}{n}
\]
does converge to $0$. It alternates, but the jumps shrink.

The near miss teaches that oscillation alone does not prevent convergence. Persistent oscillation does.

\subsection{How to present a counterexample}

A counterexample should be written as a tiny proof.

Bad:

\begin{quote}
Counterexample: $-3$.
\end{quote}

Better:

\begin{quote}
The claim ``if $x^2>4$, then $x>2$'' is false. Take $x=-3$. Then $x^2=9>4$, so the hypothesis is true, but $x>2$ is false. Therefore the implication is false.
\end{quote}

This format removes doubt. It verifies membership in the domain, checks the hypothesis, and checks the failure of the conclusion.

\begin{formalbox}
Counterexample template:

\begin{quote}
Take [object]. It belongs to [domain]. It satisfies [hypothesis] because [reason]. It fails [conclusion] because [reason]. Therefore the statement is false.
\end{quote}
\end{formalbox}

\subsection{Examples are not random}

At first, finding examples may feel like guessing. Over time, it becomes structured.

To find an injective but not surjective function, you want no merging but some missing codomain element. So shift the natural numbers:
\[
n\mapsto n+1.
\]

To find a surjective but not injective function, you want full coverage but merging. So collapse a larger domain onto a smaller codomain:
\[
\{1,2\}\to\{a\}.
\]

To find a continuous function that is not differentiable, you want no jump but a sharp corner:
\[
f(x)=|x|.
\]

To find a bounded sequence that does not converge, you want it trapped but not settling:
\[
a_n=(-1)^n.
\]

To find two different sets $B,C$ with $A\cup B=A\cup C$, make $A$ large enough to hide the difference.

This is design, not guessing.

\section{How Mathematicians Choose a Proof Strategy}

Beginners often ask, ``How was I supposed to know to try that?'' That is one of
the right questions. Proof discovery can look magical when books only show the
finished proof. In practice, mathematicians try strategies because the logical
shape of the claim suggests them.

\begin{marginnotebox}
\motifstrategy\ Strategy fork. When you see this mark, stop reading for a
moment and ask which proof path the target suggests.
\end{marginnotebox}

\subsection{Read the target first}

The conclusion usually tells you the outer proof shape.

\begin{longtable}{p{0.25\linewidth}p{0.35\linewidth}p{0.3\linewidth}}
\toprule
Target shape & Natural strategy & First move \\
\midrule
$\forall x\in A,\ P(x)$ & arbitrary object & Let $x\in A$ be arbitrary \\
$\exists x\in A,\ P(x)$ & construction & Choose or construct a witness \\
$P\Rightarrow Q$ & direct or contrapositive & Suppose $P$; or suppose $\neg Q$ \\
$P\Leftrightarrow Q$ & two directions & Prove $P\Rightarrow Q$ and $Q\Rightarrow P$ \\
$A=B$ for sets & two inclusions & Prove $A\subseteq B$ and $B\subseteq A$ \\
``unique'' & two-candidate argument & Suppose two objects satisfy the property \\
Statement over $\N$ & induction & Prove base case and step \\
Nonexistence & contradiction or universal negation & Assume existence and derive conflict \\
\bottomrule
\end{longtable}

\subsection{When contradiction is useful}

Contradiction is useful when denying the conclusion gives concrete data. For
example, to prove there is no smallest positive real number, assuming there is
one gives an object $r>0$ with a strong property. You can then construct
$r/2$, which violates that property.

Contradiction is also useful when the conclusion is negative:

\[
\text{There is no object with property }P.
\]

Assume such an object exists, then force an impossibility.

\begin{trapbox}
False intuition: contradiction is the strongest proof method, so use it
whenever stuck. Better intuition: contradiction is useful when the negation of
the target gives usable structure. If direct proof gives a clean path, use it.
\end{trapbox}

\subsection{How experts discover proofs}

Experts often work backward in scratch:

\begin{enumerate}
\item What do I need at the end?
\item Which definition would close the target?
\item What data would be enough to close that definition?
\item Do the hypotheses provide that data?
\item If not, can I prove a bridge lemma?
\end{enumerate}

Then they write the final proof forward.

\begin{labbox}
Discovery drill. For each claim, write backward scratch work before writing the
proof.
\begin{enumerate}
\item If $a\mid b$ and $a\mid c$, then $a\mid(3b-5c)$.
\item If $f$ and $g$ are injective, then $g\circ f$ is injective.
\item If $U,V\subseteq\R$ are open, then $U\cap V$ is open.
\item If $a_n\to0$, then $|a_n|\to0$.
\end{enumerate}
\end{labbox}

\subsection{Intermediate states}

Many proofs become easier when you name the missing middle. For composition of
injective functions, the final target is $a_1=a_2$. The assumption gives
$g(f(a_1))=g(f(a_2))$. The missing middle is:
\[
f(a_1)=f(a_2).
\]
Once you see that intermediate state, the proof becomes a two-step chain.

\begin{focusbox}
When stuck, ask: what statement would put me halfway between the assumptions
and the target?
\end{focusbox}

\subsection{A Recovery Page: Tiny Proofs, No Philosophy}

After a dense strategy discussion, return to small proofs. They are the scales
and arpeggios of proof writing.

\begin{exercise}
Prove that if $n$ is even, then $n+2$ is even.
\end{exercise}

\begin{proof}
Suppose $n$ is even. Then $n=2k$ for some $k\in\Z$. Thus
\[
n+2=2k+2=2(k+1).
\]
Since $k+1\in\Z$, $n+2$ is even.
\end{proof}

\begin{exercise}
Prove that if $A\subseteq B$ and $x\in A$, then $x\in B$.
\end{exercise}

\begin{proof}
This is exactly the definition of subset. Since $A\subseteq B$, every element
of $A$ is an element of $B$. Since $x\in A$, it follows that $x\in B$.
\end{proof}

\begin{recoverybox}
When the book gets conceptually heavy, do two tiny proofs. Small correct moves
restore trust in the machinery.
\end{recoverybox}

\subsection{Walking questions}

\begin{walkbox}
For every new definition, build the following list:

\begin{enumerate}
\item one tiny example,
\item one large or infinite example,
\item one non-example,
\item one near miss,
\item one boundary case,
\item one theorem where the definition is useful,
\item one false statement the definition helps refute.
\end{enumerate}
\end{walkbox}

\subsection{Exercises}

\begin{enumerate}
\item Give three examples and three non-examples of even integers. For each non-example, explain why the definition fails.
\item Give an example of two irrational numbers whose sum is rational.
\item Give an example of a rational number and an irrational number whose sum is irrational.
\item Find a counterexample: if $x^2=y^2$, then $x=y$.
\item Find a counterexample: if $x<y$, then $x^2<y^2$.
\item Find a counterexample: every function $\R\to\R$ is injective.
\item Find a counterexample: every injective function $\N\to\N$ is surjective.
\item Find a counterexample: every surjective function is injective.
\item Find a counterexample: every bounded sequence converges.
\item Find a counterexample: every continuous function is differentiable.
\item Give a near miss for the definition of prime number and explain exactly what fails.
\item Give a near miss for linear independence in $\R^2$ or $\R^3$.
\item For the theorem ``a continuous function on $[a,b]$ attains a maximum,'' give one counterexample for each missing hypothesis: continuity, closed interval, bounded interval.
\item For each false statement above, rewrite it into a nearby true statement.
\item Choose a definition from a later course you care about and build a mini zoo: examples, non-examples, counterexamples, and boundary cases.
\end{enumerate}
